\documentclass[10pt]{amsart}
\usepackage[a4paper,margin=23mm]{geometry}
\usepackage{amsmath,amssymb,amsthm,mathtools,hyperref,nicefrac}
\newtheorem{prop}{Proposition}[section]
\newtheorem{lemm}[prop]{Lemma}
\newtheorem{coro}[prop]{Corollary}
\newtheorem*{rema*}{Remark}
\newcommand{\eg}{e.g.\ }
\emergencystretch=3em
\title{Monotonicity of bar-weight under coprime bar-core reduction}
\author{Dean Yates}
\date{}
\begin{document}
\maketitle
\noindent\textbf{Source.} Original Proposition 5.1, A generalisation of bar-core partitions, Algebraic Combinatorics 5 (2022), 667--698, DOI 10.5802/alco.231. For coprime odd $p,q\geq3$, all original bar-core/weight/quotient and level-group-action definitions and required local lemmas, generalized-weight comparison and orbit classification proofs precede the complete weight-comparison proof and original equality criterion. Numerical examples and their abacus/Young-diagram illustrations are ancillary and omitted, with original proposition numbering preserved. No original mathematical argument is corrected or supplemented.
\section{Introduction}
 
The study of the set of core partitions, those partitions whose Young diagrams are without $s$-hooks for some natural number $s$, has revealed a great deal about the representation theory of the symmetric groups, since two irreducible characters of $S_n$ are in the same $s$-block if and only if the partitions labelling them have the same $s$-core. The motivation behind the study of \textit{bar partitions} is their correspondence to \textit{projective}, or \textit{spin representations} of the symmetric group \cite{HH}. The irreducible spin representations of $S_n$ corresponding to bar partitions $\lambda$ and $\mu$ lie in the same $p$-block if and only if they have the same $p$\textit{-bar-core}.  

The purpose of this paper is to establish analogues for bar partitions, i.e. partitions with distinct parts (or simply finite subsets of $\mathbb{N}$), of the results in Fayers' `A generalisation of core partitions' \cite{F1}, which includes a generalisation of Olsson's result that shows the $s$-core of a $t$-core is again a $t$-core \cite{O2}. Although the definitions differ, the ideas in \cite{F1} translate very well to the notion of $p$-bar-cores, and this paper includes a generalisation of a further result established by Olsson in \cite{O2}: that the $q$-bar-core of a $p$-bar-core is again a $p$-bar-core. 

We will begin with a few definitions, which will seem very familiar to those acquainted with the representation theory of the symmetric group, that have been adapted to suit our purposes. Using some basic results, we consider an action of $\mathfrak{W}_p$, the Weyl group of type $\tilde{C}_{(p-1)/2}$, on the set of bar partitions $\overline{\mathcal{P}}$, and discover some interesting symmetry. We then consider the problem of constructing the smallest bar partition with a given $p$-bar-core and $q$-bar-core, for coprime odd $p,q\geq3$. We finish by investigating the orbits of the Yin and Yang partitions \cite{BO} under the group action of $\mathfrak{W}_p\times\mathfrak{W}_q$. 

\section{Definitions}
 
A \textbf{bar partition} $\lambda\in\overline{\mathcal{P}}$ is a decreasing sequence of distinct positive integers (often referred to as a 2-regular or strict partition). For odd integers $p\geq3$, the $p$\textbf{-runner abacus} \cite{F2} has $p$ infinite vertical runners numbered from left to right $\nicefrac{(p+1)}{2},\nicefrac{(p+3)}{2},\dots,p-1,0,1,\dots,\nicefrac{(p-1)}{2}$, with the positions on runner $i$ labelled with the integers with $p$ residue $i$, increasing down the runner, so that position $x+1$ appears directly to the right of position $x$ (for $x\not\equiv\nicefrac{(p-1)}{2}$ mod $p$). We obtain a visual representation of $\lambda$ on the $p$-runner abacus by placing a bead on position $x$ for each $x\in\lambda$ and each integer $x<0$ such that $-x\not\in\lambda$; position 0 remains empty. 

This differs from the way that Olsson, for example, represents bar partitions and bar-cores, but this $p$-runner abacus will be more useful for our purposes. 

\stepcounter{prop}

Let $\mathcal{A}(\lambda)$ denote the set containing all integers that label positions occupied by beads in the bead configuration for $\lambda\in\overline{\mathcal{P}}$ on the $p$-runner abacus: 
$$x\in\mathcal{A}(\lambda)\Leftrightarrow\begin{cases}x\in\lambda,&x>0\text{, or}\\-x\not\in\lambda,&x<0.\end{cases}$$
Note that this is independent from the choice of $p$. 
\newline
For odd integers $p\geq3$, removing a $p$\textbf{-bar} from $\lambda\in\overline{\mathcal{P}}$ means either \newline
(i) removing $x\in\lambda$ such that $0\leq x-p\not\in\lambda$, and replacing $x$ with $x-p$ if $x\neq p$; or \newline
(ii) removing two parts $x,p-x\in\lambda$ (where $0<x<p$).
\newline
 ($p$ must be odd because of the incompatible possibility that a bar partition could have a $2p$-bar but not a $p$-bar, \eg $p=4$ and the partition $(6,2)$.) \newline
In terms of the abacus, removing a $p$-bar from $\lambda$ corresponds to moving a bead at position $x$ to position $x-p$ (replacing $x\in\mathcal{A}(\lambda)$ with $x-p$), then moving the bead at position $p-x$ to position $-x$ (replacing $p-x\in\mathcal{A}(\lambda)$ with $-x$). 


When it is not possible to remove any $p$-bars from $\lambda$, i.e. when $x-p\in\mathcal{A}(\lambda)$ for all $x\in\mathcal{A}(\lambda)$, we say that $\lambda$ is a $p$\textbf{-bar-core}, and we denote the set of $p$-bar-cores by $\overline{C}_p$. 
\newline
Since removing a $p$-bar always corresponds to moving beads up on their runners to unoccupied positions, we have reached the bead configuration of a $p$-bar-core when all beads are moved up their runners as far as possible. 
The order in which these moves are made is irrelevant; we always end up at the same bead configuration. 
Hence we may define \textbf{the} $p$\textbf{-bar-core} of a bar partition $\lambda$, which we denote by $\overline{\lambda}_p$. 
The number of $p$-bars which can be successively removed from $\lambda$ is the $p$\textbf{-bar-weight} of $\lambda$, and we denote this quantity by $\overline{\text{wt}}_p(\lambda)$; denoting by $|\mu|$ the sum of the parts of the bar partition $\mu$, 
$$\overline{\text{wt}}_p(\lambda)\coloneqq\tfrac{|\lambda|-|\overline{\lambda}_p|}{p}.$$ 
The number of bead moves needed to reach the bead configuration for $\overline{\lambda}_p$ from the bead configuration for $\lambda$ is equal to twice the $p$-bar-weight of $\lambda$, because removing a $p$-bar corresponds to two moves of the beads. 
The $p$-bar-weight of $\lambda$ is therefore equal to half the number of pairs $(x,a)\in\mathcal{A}(\lambda)\times\mathbb{N}$ such that $x-ap\not\in\mathcal{A}(\lambda)$. 

\stepcounter{prop}

We are now equipped with tools analogous to those James used in his seminal book on the representation theory of the symmetric group via the combinatorics of (not necessarily strict) partitions \cite{JK}. For the benefit of readers unfamiliar with James' work, we outline the theory of rim-hooks and cores here.

We may visually represent a partition $\alpha=(\alpha_1,\alpha_2,\dots,\alpha_r)$, i.e. a decreasing sequence of (not necessarily distinct) positive integers $\alpha_1\geq\alpha_2\geq\dots\geq\alpha_r$, by its \textbf{Young diagram} $[\alpha]$, which has $\alpha_i$  nodes in the $i^\text{th}$ row, for each $i\in\{1,\dots,r\}$, with each row starting in the first column. The \textbf{hook-length} $h_{i,j}$ of the $(i,j)$-node, in the $i^\text{th}$ row and $j^\text{th}$ column of $[\alpha]$, is found by adding the number of $(k,j)$-nodes with $k\geq i$ to the number of $(i,l)$-nodes with $l>j$. We refer to the $(i,j)$-nodes with $(i+1,j+1)\not\in[\alpha]$ as the \textbf{rim} of $[\alpha]$. 
The $h_{i,j}$ pairwise adjacent nodes along the rim of $[\alpha]$ from the lowest node in the $j^\text{th}$ column, i.e. the $(k,j)$-node with $k$ maximal, to the $(i,\alpha_i)$-node are collectively called a \textbf{rim} $h_{i,j}$-\textbf{hook}. Whenever a diagram $[\alpha]$ has an $(i,j)$-node with hook-length $s\coloneqq h_{i,j}$, we may remove a rim $s$-hook from $[\alpha]$ to obtain the Young diagram of a different partition. If instead $[\alpha]$ has no nodes with hook-length $s$, then we say that the partition $\alpha$ is an $s$-\textbf{core}. 

\stepcounter{prop}

Adopting the convention that $\alpha_i=0$ for each $i$ greater than some fixed $r\in\mathbb{N}$, the strictly decreasing sequence of integers $\alpha_1-1+k>\alpha_2-2+k>\dots$, for some $k\in\mathbb{Z}$, is called a \textbf{beta}-\textbf{set} for the partition $\alpha=(\alpha_1,\alpha_2,\dots)$, and is denoted by $\mathcal{B}^\alpha_k$. 


James' $s$\textbf{-abacus} has $s$ runners extending infinitely in both directions, with the leftmost runner labelled by multiples of $s$, and the position directly to the right of $i$ labelled by $i+1$. A bead configuration is associated with a partition $\alpha$ via the beta-set $\mathcal{B}^\alpha\coloneqq\mathcal{B}^\alpha_0$ by placing a bead at the position labelled by $\alpha_i-i$ for each $i\in\mathbb{N}$. 

Removing a rim $s$-hook from $[\alpha]$ then corresponds to removing an element $x\in\mathcal{B}^\alpha$ such that $x-s\not\in\mathcal{B}^\alpha$, and replacing $x$ with $x-s$. Thus we obtain the bead configuration for an $s$-core by moving the beads in the configuration for $\alpha$ on the $s$-abacus up their runners as far as possible. Since the order in which we move the beads is irrelevant, there is only one $s$-core which can be obtained from a partition $\alpha$ by removing rim $s$-hooks, and we denote \textbf{the} $s$\textbf{-core} of $\alpha$ by $\tilde{\alpha}_s$. The number of moves needed to reach the bead configuration of $\tilde{\alpha}_s$ from the configuration of $\alpha$, or equivalently, the number of rim $s$-hooks which can be removed from the diagram $[\alpha]$, is the $s$-\textbf{weight} of $\alpha$; we denote this quantity by $\text{wt}_s(\alpha)$. 

The $s$-\textbf{quotient} of $\alpha$ is the $s$-tuple of partitions corresponding to the bead configuration of each runner of the $s$-abacus as $s$ separate 1-abaci. Each partition $\alpha$ is uniquely determined by its $s$-core $\tilde{\alpha}_s$ and its $s$-quotient. 

\stepcounter{prop}

Removing rim $s$-hooks has a strong connection with the modular representation theory of the symmetric group, as the two ordinary irreducible representations corresponding to the partitions $\alpha$ and $\beta$ belong to the same $s$-\textit{block} of $s$-modular irreducible constituents if and only if $\tilde{\alpha}_s=\tilde{\beta}_s$. This important result was first conjectured by Nakayama, and should be referred to as the Brauer-Robinson Theorem after those who first proved it in 1947.

Now that we have seen how the combinatorics of bar partitions is related to James' work, we introduce a useful way to encode bar partitions, just as the $s$-core and $s$-quotient encode partitions \cite{JK}.

Define the $p$\textbf{-set} \cite{F3} of a bar partition $\lambda$ to be the set $\{\Delta_{i\bmod p}\lambda|i\equiv0,1,\dots,p-1\}$, where $\Delta_{i\bmod p}\lambda$ is the smallest integer $x\equiv i$ modulo $p$ such that $x\not\in\mathcal{A}(\overline{\lambda}_p)$. 
Since $x\in\mathcal{A}(\overline{\lambda}_p)\Leftrightarrow-x\not\in\mathcal{A}(\overline{\lambda}_p)$, for any bar partition $\lambda$ and $k\not\equiv0$ (mod $p$) we have $\Delta_{k\bmod p}\lambda+\Delta_{-k\bmod p}\lambda=p$, so all of the elements in the $p$-set of any bar partition sum to $\nicefrac{p(p-1)}{2}$. 

Next we will introduce the $p$\textbf{-quotient} of $\lambda\in\overline{\mathcal{P}}$ \cite{O1}, which is the $p$-tuple 
$$\mathcal{Q}_p(\lambda)\coloneqq(\lambda^{(0\bmod p)},\dots,\lambda^{(p-1\bmod p)}).$$  
(We will drop the `mod $p$' in our notation for both the $p$-set and $p$-quotient when it is clear which $p$ we are referring to.) 

The parts of the bar partition $\lambda^{(0\bmod p)}$ are the elements of the set $\{\nicefrac{x}{p}|x\in\lambda,x\in p\mathbb{Z}\}$. 
\newline
For $j\not\equiv0$ (mod $p$), the $i^\text{th}$ part of the (not necessarily strict) partition $\lambda^{(j\bmod p)}$ is equal to the number of empty spaces above the $i^\text{th}$ lowest bead on runner $j$ in the bead configuration for $\lambda$ on the $p$-runner abacus. 

It follows from the definition of the $p$-runner abacus that for each $j\not\equiv0$ (mod $p$), the partition $\lambda^{(-j\bmod p)}$ is equal to $(\lambda^{(j\bmod p)})'$, the \textit{conjugate} of the partition $\lambda^{(j\bmod p)}$, the parts of which are the lengths of the columns in the Young diagram for $\lambda^{(j\bmod p)}$. 

Each $p$-set corresponds to a unique $p$-bar-core, and Olsson proved that every bar partition is uniquely determined by its $p$-bar-core and $p$-quotient \cite[Proposition 2.2]{O1}. 

\stepcounter{prop}

\section{Properties of the $p$-quotient}
 
Later we will utilise the following properties of $\mathcal{Q}_p(\lambda)$. 
\begin{lemm}\label{lem3}
Suppose $\lambda\in\overline{\mathcal{P}}$ and $c,p$ are odd positive integers, with $p\geq 3$. 
\begin{enumerate}
\item $\overline{\text{wt}}_{cp}(\lambda)=\overline{\text{wt}}_{c}(\lambda^{(0\bmod p)})+\tfrac{1}{2}\sum_{j=1}^{p-1}\text{wt}_c(\lambda^{(j\bmod p)})$; in particular, $\lambda$ is a $cp$-bar-core if and only if $\lambda^{(0\bmod p)}$ is a $c$-bar-core and every other component of the $p$-quotient of $\lambda$ is a $c$-core. 
\item ${\overline{(\lambda^{(0\bmod p)}})}_c=(\overline{\lambda}_{cp})^{(0\bmod p)}$, and for $j\not\equiv0\:(\text{mod }p)$, the $c$-core of $\lambda^{(j\bmod p)}$ is equal to $(\overline{\lambda}_{cp})^{(j\bmod p)}$.
\end{enumerate}
\end{lemm}
\begin{proof}
Removing a $cp$-bar from $\lambda$ means replacing an element $x\in\lambda$ with $x-cp$ (when $x-cp\not\in\mathcal{A}(\lambda)$), then replacing $cp-x\in\mathcal{A}(\lambda)$ with $-x$. 

If $x\equiv0$ (mod $p$), then this is equivalent to replacing $\nicefrac{x}{p}\in\mathcal{A}(\lambda^{(0)})$ with $\nicefrac{(x-cp)}{p}=\nicefrac{x}{p}-c$, and replacing $\nicefrac{(cp-x)}{p}\in\mathcal{A}(\lambda^{(0)})$ with $\nicefrac{-x}{p}=\nicefrac{(cp-x)}{p}-c$, i.e. removing a $c$-bar from $\lambda^{(0)}$. 

If $x\equiv j$ (mod $p$), then this is equivalent to replacing $\nicefrac{(x-j)}{p}\in\mathcal{B}^{\lambda^{(j)}}_r$ with $\nicefrac{((x-cp)-j)}{p}=\nicefrac{(x-j)}{p}-c$, and then replacing $\nicefrac{((cp-x)-(p-j))}{p}\in\mathcal{B}^{\lambda^{(p-j)}}_r$ with $\nicefrac{(-x-(p-j))}{p}=\nicefrac{((cp-x)-(p-j))}{p}-c$, thus removing a rim $c$-hook from both $\lambda^{(j)}$ and $\lambda^{(p-j)}$.
\end{proof}

\begin{rema*}
Note that since $\lambda^{(-j\bmod p)}=(\lambda^{(j\bmod p)})'$ for $j\not\equiv0$ (mod $p$), we can rewrite the first part of the previous lemma as 
$$\overline{\text{wt}}_{cp}(\lambda)=\overline{\text{wt}}_{c}(\lambda^{(0\bmod p)})+\sum_{j=1}^{\tfrac{p-1}{2}}\text{wt}_c(\lambda^{(j\bmod p)}).$$ 
\end{rema*}

\section{A level $q$ group action on bar partitions}
 
Now we will consider an action of $\mathfrak{W}_p$, the affine Coxeter group of type $\tilde{C}_{\nicefrac{(p-1)}{2}}$, with generators $\delta_0, \dots, \delta_{\nicefrac{(p-1)}{2}}$ and relations 
\begin{align*}\delta_i^2&=1&&\text{for }0\leq{i}\leq\tfrac{p-1}{2},\\
\delta_i\delta_j&=\delta_j\delta_i&&\text{for }0\leq{i}<j-1\leq\tfrac{p-3}{2},\\
\delta_i\delta_{i+1}\delta_i&=\delta_{i+1}\delta_i\delta_{i+1}&&\text{for }1\leq{i}\leq\tfrac{p-5}{2},\\
\delta_0\delta_1\delta_0\delta_1&=\delta_1\delta_0\delta_1\delta_0&&\text{if }p>3,\\
\delta_{\nicefrac{(p-3)}{2}}\delta_{\nicefrac{(p-1)}{2}}\delta_{\nicefrac{(p-3)}{2}}\delta_{\nicefrac{(p-1)}{2}}&=\delta_{\nicefrac{(p-1)}{2}}\delta_{\nicefrac{(p-3)}{2}}\delta_{\nicefrac{(p-1)}{2}}\delta_{\nicefrac{(p-3)}{2}}&&\text{if }p>3.\end{align*}
For coprime odd integers $p,q\geq3$, we define a \textit{level} $q$ \textit{action} of $\mathfrak{W}_p$ on $\mathbb{Z}$ \cite{F3}: 
\begin{align*}\delta_0x&=\begin{cases}x-2q&\text{if }x\equiv q\:(\text{mod }p),\\
x+2q&\text{if }x\equiv-q\:(\text{mod }p),\\
x&\text{otherwise;}\end{cases}\\ 
\delta_ix&=\begin{cases}x-q&\text{if }x\equiv(i+1)q,-iq\:(\text{mod }p),\\
x+q&\text{if }x\equiv iq,-(i+1)q\:(\text{mod }p),\\
x&\text{otherwise,}\end{cases}&&\text{for }1\leq i\leq\tfrac{p-1}{2}.\end{align*}
For the rest of this paper, we will assume that $p$ and $q$ are coprime odd integers no less than 3. 
\begin{lemm}
The above defines a group action of $\mathfrak{W}_p$ on $\mathbb{Z}$, and this can be extended to an action on $\overline{\mathcal{P}}$.
\end{lemm}
\begin{proof}
We must have $p>3$ for the fourth and fifth relations of the generators of $\mathfrak{W}_p$ to hold for the group action. 
The first relation $\delta_i^2=1$ is clear for all $i$. Moreover, the generators $\delta_i,\delta_j$ commute when $0\leq i<j-1\leq \nicefrac{(p-3)}{2}$ because they act on distinct congruence classes of integers modulo $p$. For the third relation, when $1\leq i\leq\nicefrac{(p-5)}{2}$ and $x\in\mathbb{Z}$ we have 
\begin{align*}
\delta_i\delta_{i+1}\delta_ix&=\left\{\begin{array}{l|l}\delta_i(x-2q)&x\equiv-iq\\
\delta_i(x+2q)&x\equiv iq\\
\delta_i(x-q)&x\equiv(i+1)q,(i+2)q\\
\delta_i(x+q)&x\equiv-(i+1)q,-(i+2)q\\
\delta_ix&\text{otherwise}\end{array}\right\}&&=\left\{\begin{array}{l|l}x-2q&x\equiv(i+2)q,-iq\\
x+2q&x\equiv iq,-(i+2)q\\
x&\text{otherwise}\end{array}\right.\\
&=\left\{\begin{array}{l|l}\delta_{i+1}(x-2q)&x\equiv(i+2)q\\
\delta_{i+1}(x+2q)&x\equiv-(i+2)q\\
\delta_{i+1}(x-q)&x\equiv-iq,-(i+1)q\\
\delta_{i+1}(x+q)&x\equiv iq,(i+1)q\\
\delta_{i+1}x&\text{otherwise}\end{array}\right\}&&=\delta_{i+1}\delta_i\delta_{i+1}x.\end{align*}
For the fourth and fifth relations, assuming $p>3$, 
\begin{align*}(\delta_0\delta_1)^2x&=\left\{\begin{array}{l|l}\delta_0\delta_1(x-3q)&x\equiv2q\\
\delta_0\delta_1(x+3q)&x\equiv-2q\\
\delta_0\delta_1(x-q)&x\equiv-q\\
\delta_0\delta_1(x+q)&x\equiv q\\
\delta_0\delta_1x&\text{otherwise}\end{array}\right\}&&=\left\{\begin{array}{l|l}x-4q&x\equiv2q\\
x+4q&x\equiv-2q\\
x-2q&x\equiv q\\
x+2q&x\equiv-q\\
x&\text{otherwise}\end{array}\right.\\
&=\left\{\begin{array}{l|l}\delta_1\delta_0(x-3q)&x\equiv q\\
\delta_1\delta_0(x+3q)&x\equiv-q\\
\delta_1\delta_0(x-q)&x\equiv2q\\
\delta_1\delta_0(x+q)&x\equiv-2q\\
\delta_1\delta_0x&\text{otherwise}\end{array}\right\}&&=(\delta_1\delta_0)^2x,\\
(\delta_{\nicefrac{(p-3)}{2}}\delta_{\nicefrac{(p-1)}{2}})^2x&=\left\{\begin{array}{l|l}\delta_{\nicefrac{(p-3)}{2}}\delta_{\nicefrac{(p-1)}{2}}(x-2q)&x\equiv\tfrac{(p+1)q}{2}\\
\delta_{\nicefrac{(p-3)}{2}}\delta_{\nicefrac{(p-1)}{2}}(x+2q)&x\equiv\tfrac{(p-1)q}{2}\\
\delta_{\nicefrac{(p-3)}{2}}\delta_{\nicefrac{(p-1)}{2}}(x-q)&x\equiv\tfrac{(p+3)q}{2}\\
\delta_{\nicefrac{(p-3)}{2}}\delta_{\nicefrac{(p-1)}{2}}(x+q)&x\equiv\tfrac{(p-3)q}{2}\\
\delta_{\nicefrac{(p-3)}{2}}\delta_{\nicefrac{(p-1}{2}}x&\text{otherwise}\end{array}\right\}&&=\left\{\begin{array}{l|l}x-3q&x\equiv\tfrac{(p+3)q}{2}\\
x+3q&x\equiv\tfrac{(p-3)q}{2}\\
x-q&x\equiv\tfrac{(p+1)q}{2}\\
x+q&x\equiv\tfrac{(p-1)q}{2}\\
x&\text{otherwise}\end{array}\right.\\
&=\left\{\begin{array}{l|l}\delta_{\nicefrac{(p-1)}{2}}\delta_{\nicefrac{(p-3)}{2}}(x-2q)&x\equiv\tfrac{(p+3)q}{2}\\
\delta_{\nicefrac{(p-1)}{2}}\delta_{\nicefrac{(p-3)}{2}}(x+2q)&x\equiv\tfrac{(p-3)q}{2}\\
\delta_{\nicefrac{(p-1)}{2}}\delta_{\nicefrac{(p-3)}{2}}(x-q)&x\equiv\tfrac{(p-1)q}{2}\\
\delta_{\nicefrac{(p-1)}{2}}\delta_{\nicefrac{(p-3)}{2}}(x+q)&x\equiv\tfrac{(p+1)q}{2}\\
\delta_{\nicefrac{(p-1)}{2}}\delta_{\nicefrac{(p-3)}{2}}x&\text{otherwise}\end{array}\right\}&&=(\delta_{\nicefrac{(p-1)}{2}}\delta_{\nicefrac{(p-3)}{2}})^2x.
\end{align*}

If $X$ is a subset of $\mathbb{Z}\backslash\{0\}$ that is bounded above, and its complement in $\mathbb{Z}$ is bounded below, then it is easy to see that the same is true for $aX\coloneqq\{ax|x\in X\}$, for any $a\in\mathfrak{W}_p$. Moreover, when $x\in X\Leftrightarrow-x\not\in X$, for all $x\in\mathbb{Z}\backslash\{0\}$, then the set $aX$ also satisfies this rule. 
Hence, this action can be extended to an action on bar partitions $\lambda$ by defining $\delta_i\lambda$ to be the bar partition with $\mathcal{A}(\delta_i\lambda)=\delta_i\mathcal{A}(\lambda)$.
\end{proof}

\stepcounter{prop}

We now give some invariants of the level $q$ action of $\mathfrak{W}_p$ which we will later use to give an explicit criterion for when two bar partitions lie in the same orbit under the level $q$ action, which we refer to as a \textbf{level} $q$ \textbf{orbit}.

\begin{lemm}\label{lem4} 
Suppose $\lambda\in\overline{\mathcal{P}}$ and $a\in\mathfrak{W}_p$, and define $a\lambda$ using the level $q$ action. 
\begin{enumerate}
\item $\overline{(a\lambda)}_q=\overline{\lambda}_q$; 
\item $\mathcal{Q}_p(a\lambda)$ is the same as $\mathcal{Q}_p(\lambda)$ with the components reordered; 
\item $\overline{\text{wt}}_p(a\lambda)=\overline{\text{wt}}_p(\lambda)$; 
\item $\overline{(a\lambda)}_p=a(\overline{\lambda}_p)$.
\end{enumerate}
\end{lemm}
\begin{proof}
The relations occurring in all four parts are transitive, so we need only prove them in the case where $a$ is simply a generator $\delta_i$ of $\mathfrak{W}_p$. 

\smallskip
\noindent
(1) An element $x\in\mathcal{A}(\lambda)$ is fixed by the level $q$ action of $\delta_0$ if and only if 
$$x\not\equiv q,-q\:(\text{mod }p);\:x\equiv q\:(\text{mod }p),\:x-2q\in\mathcal{A}(\lambda);\text{ or }x\equiv-q\:(\text{mod }p),\:x+2q\in\mathcal{A}(\lambda).$$ 
If $x\in\mathcal{A}(\lambda)$, $x\equiv q$ (mod $p$) and $x-2q\not\in\mathcal{A}(\lambda)$, then $\delta_0$ sends $x$ to $x-2q$, and sends $2q-x\in\mathcal{A}(\lambda)$ to $-x$. If $x\in\mathcal{A}(\lambda)$, $x\equiv-q$ (mod $p$) and $x+2q\not\in\mathcal{A}(\lambda)$, then $\delta_0$ sends $x$ to $x+2q$, and sends $-x-2q\in\mathcal{A}(\lambda)$ to $-x$. Thus the action of $\delta_0$ on $\mathcal{A}(\lambda)$ corresponds to removing $2q$-bars from or adding $2q$-bars to $\lambda$. 

For $i\in\{1,\dots,\nicefrac{(p-1)}{2}\}$, the level $q$ action of $\delta_i$ fixes $x\in\mathcal{A}(\lambda)$ if and only if 
\begin{align*}x\not\equiv(i+1)q,-iq,iq,-(i+1)q\:(\text{mod }p);\:x\equiv(i+1)q,-iq\:(\text{mod }p),\:x-q\in\mathcal{A}(\lambda);\\
\text{or }x\equiv iq,-(i+1)q\:(\text{mod }p),\:x+q\in\mathcal{A}(\lambda).\end{align*} 
If $x\in\mathcal{A}(\lambda)$, $x\equiv(i+1)q$ or $-iq$ (mod $m$) and $x-q\not\in\mathcal{A}(\lambda)$, then $\delta_i$ sends $x$ to $x-q$, and sends $q-x\in\mathcal{A}(\lambda)$ to $-x$. If $x\in\mathcal{A}(\lambda)$, $x\equiv iq$ or $-(i+1)q$ (mod $p$) and $x+q\not\in\mathcal{A}(\lambda)$, then $\delta_i$ sends $x$ to $x+q$, and sends $-x-q\in\mathcal{A}(\lambda)$ to $-x$. Thus the action of $\delta_i$ on $\mathcal{A}(\lambda)$ corresponds to removing $q$-bars from or adding $q$-bars to $\lambda$. 

Moreover, since there are only finitely many $x\in\mathcal{A}(\lambda)$ such that at least one of $x-2q$, $x+2q$, $x-q$ or $x+q$ is not in $\mathcal{A}(\lambda)$, we obtain $\mathcal{A}(\delta_i\lambda)\coloneqq\delta_i\mathcal{A}(\lambda)$ from $\mathcal{A}(\lambda)$ via a finite number of changes. Hence $\overline{(\delta_i\lambda)}_q=\overline{\lambda}_q$. 

\smallskip
\noindent
(2) When $x\not\equiv\pm q$ (mod $p$), we have $x\in\mathcal{A}(\delta_0\lambda)\Leftrightarrow x\in\mathcal{A}(\lambda)$, so $(\delta_0\lambda)^{(j)}=\lambda^{(j)}$ for all $j\not\equiv\pm q$ (mod $p$). If instead $x\equiv q$ (mod $p$), then $x\in\mathcal{A}(\delta_0\lambda)\Leftrightarrow x+2q\in\mathcal{A}(\lambda)$, or if $x\equiv-q$ (mod $p$), then $x\in\mathcal{A}(\delta_0\lambda)\Leftrightarrow x-2q\in\mathcal{A}(\lambda)$; hence $(\delta_0\lambda)^{(q)}=\lambda^{(-q)}=(\lambda^{(q)})'$ and $(\delta_0\lambda)^{(-q)}=\lambda^{(q)}=(\lambda^{(-q)})'$. 

If $i\in\{1,\dots,\nicefrac{(p-1)}{2}\}$ and $x\not\equiv(i+1)q,-iq,iq,-(i+1)q$ (mod $p$), then $x\in\mathcal{A}(\delta_i\lambda)\Leftrightarrow x\in\mathcal{A}(\lambda)$, so $(\delta_i\lambda)^{(j)}=\lambda^{(j)}$ for all $j\not\equiv(i+1)q,-iq,iq,-(i+1)q$ (mod $p$). If instead $x\equiv(i+1)q$ or $-iq$ (mod $p$), then $x\in\mathcal{A}(\delta_i\lambda)\Leftrightarrow x+q\in\mathcal{A}(\lambda)$, and if $x\equiv iq$ or $-(i+1)q$ (mod $p$), then $x\in\mathcal{A}(\delta_i\lambda)\Leftrightarrow x-q\in\mathcal{A}(\lambda)$; hence if $j\equiv(i+1)q$ or $-iq$ (mod $p$), then $(\delta_i\lambda)^{(j)}=\lambda^{(j-q)}$ and $(\delta_i\lambda)^{(-j)}=\lambda^{(j+q)}$. 

\smallskip
\noindent
(3) This follows from (2) and Lemma \ref{lem3}(1) (taking $c=1$): the bead configuration for $\delta_i\lambda$ on the $p$-runner abacus is the same as that of $\lambda$ but with the runners reordered, so the $p$-bar-weights of the two bar partitions are equal. 

\smallskip
\noindent
(4) We need to show that $\Delta_{j\bmod p}(\delta_i\lambda)=\delta_i(\Delta_{k\bmod p}\lambda)$, for each $i\in\{0,\dots,\nicefrac{(p-1)}{2}\}$ and $j\equiv\delta_ik$ (mod $p$). We suppose the contrary. 

When $i=0$, we may assume that $j\equiv\pm q\equiv-k$ (mod $p$), as otherwise $j=k$ and 
$$\Delta_j(\delta_0\lambda)=\Delta_j\lambda=\Delta_k\lambda=\delta_0(\Delta_k\lambda).$$ 
But $$j\equiv\pm q\equiv-k\Rightarrow\Delta_j(\delta_0\lambda)=\Delta_k\lambda\pm2q=\delta_0(\Delta_k\lambda),$$ 
so we must have $i>0$, and we may also assume that $j\equiv\pm iq$ or $\pm(i+1)q$, and $j\equiv k\pm q$ (mod $p$). But then 
$$\Delta_j(\delta_i\lambda)=\Delta_{\delta_i(j)}\lambda=\delta_i(\Delta_j\lambda),$$ 
so in fact, the $p$-set of $\delta_i\lambda$ is equal to the image of the $p$-set of $\lambda$ under the action of $\delta_i$. Hence $\overline{(\delta_i\lambda)}_p=\delta_i(\overline{\lambda}_p)$, for all $i\in\{0,\dots,\nicefrac{(p-1)}{2}\}$. 
\end{proof}

Next we will give a criterion for when two bar partitions lie in the same level $q$ orbit; to this end, we will first establish a condition for two $p$-bar-cores to lie in the same level $q$ orbit.

\begin{prop}\label{prop7} 
Suppose $\lambda,\mu\in\overline{C}_p$, and that the multisets 
$$[\Delta_i\lambda\:(\text{mod }q)|i\equiv0,\dots,p-1\:(\text{mod }p)],\:[\Delta_i\mu\:(\text{mod }q)|i\equiv0,\dots,p-1\:(\text{mod }p)]$$
are equal. Then $\overline{\lambda}_q=\overline{\mu}_q$, and $\lambda$ and $\mu$ lie in the same level $q$ orbit \textup{(}of $p$-bar-cores\textup{)}. 
\end{prop}
\begin{proof}
The fact that $\lambda$ and $\mu$ have the same $q$-bar-core is established by Fayers' \cite[Proposition 4.1]{F3}, and it follows from the definition of the level $q$ action of $\mathfrak{W}_p$ on the set of $p$-bar-cores that this action preserves the $q$-bar-core of a bar partition as by definition $\delta_i$ does not change the multiset of residues modulo $q$ of the elements of the $p$-set. Therefore each orbit of the level $q$ action on $\overline{C}_p$ can contain at most one $q$-bar-core. 

In the same paper \cite{F3}, Fayers proves the following result: 
\begin{center}
Suppose $\mathcal{O}$ is a level $q$ orbit. Let $\nu$ be an element of $\mathcal{O}$ for which the sum $\sum^{p-1}_{i=0}(\Delta_{i\bmod p}\nu-\nicefrac{p}{2})^2$ is minimised. Then $\nu$ is a $q$-bar-core. 
\end{center}
This $\nu$ is uniquely defined as each level $q$ orbit contains no more than one $q$-bar-core. Thus letting $\nu$ be the $q$-bar-core of both $\lambda$ and $\mu$, it must be contained in both the level $q$ orbit containing $\lambda$ and the level $q$ orbit containing $\mu$; so these orbits coincide. 
\end{proof}

For the more general result, we define the $q$\textbf{-weighted} $p$\textbf{-quotient} of $\lambda\in\overline{\mathcal{P}}$ with $p$-set $\{\Delta_0\lambda,\dots,\Delta_{p-1}\lambda\}$ and $p$-quotient $\mathcal{Q}_p(\lambda)=(\lambda^{(0)},\dots,\lambda^{(p-1)})$ to be the multiset 
$$\mathcal{Q}^q_p(\lambda)\coloneqq[(\Delta_i\lambda\:(\text{mod }q),\lambda^{(i)})|i\equiv0,1,\dots,p-1].$$ 

\stepcounter{prop}

\begin{prop}\label{prop8} 
Suppose $\lambda,\mu\in\overline{\mathcal{P}}$. Then $\lambda$ and $\mu$ lie in the same level $q$ orbit of $\mathfrak{W}_p$ if and only if they have the same $q$-weighted $p$-quotient. 
\end{prop}
\begin{proof}
Firstly suppose that $\lambda$ and $\mu$ lie in the same level $q$ orbit; we may assume that $\mu=\delta_i\lambda$ for some $i\in\{0,\dots,\nicefrac{(p-1)}{2}\}$. Then we get $\mathcal{Q}^q_p(\lambda)=\mathcal{Q}^q_p(\mu)$ from the proof of Lemma \ref{lem4}: when $\mu=\delta_0\lambda$, we have 
$$(\Delta_j\mu,\mu^{(j)})=\begin{cases}(\Delta_{-j}\lambda+2q,\lambda^{(-j)})&\text{for }j\equiv q\text{ (mod }p),\\
(\Delta_{-j}\lambda-2q,\lambda^{(-j)})&\text{for }j\equiv-q\text{ (mod }p),\\
(\Delta_j\lambda,\lambda^{(j)})&\text{otherwise},\end{cases}$$ 
and when $\mu=\delta_i\lambda$ for some $i\in\{1,\dots,\nicefrac{(p-1)}{2}\}$, we have 
$$(\Delta_j\mu,\mu^{(j)})=\begin{cases}(\Delta_{j-q}\lambda+q,\lambda^{(j-q)})&\text{for }j\equiv(i+1)q\text{ or }-iq\text{ (mod }p),\\
(\Delta_{j+q}\lambda-q,\lambda^{(j+q)})&\text{for }j\equiv iq\text{ or }-(i+1)q\text{ (mod }p),\\
(\Delta_j\lambda,\lambda^{(j)})&\text{otherwise}.\end{cases}$$ 

For the other direction, suppose that $\lambda$ and $\mu$ share the $q$-weighted $p$-quotient $\mathcal{Q}^q_p(\lambda)=\mathcal{Q}^q_p(\mu)$. By the definition of the $p$-set, and since all components of the $p$-quotient of a $p$-bar-core are equal to the empty bar partition, the $p$-bar-cores of $\lambda$ and $\mu$ must have the same $q$-weighted $p$-quotient $\mathcal{Q}^q_p(\overline{\lambda}_p)=\mathcal{Q}^q_p(\overline{\mu}_p)$. Thus, by Proposition \ref{prop7} we may find $a,b\in\mathfrak{W}_p$ such that $a(\overline{\lambda}_p)=b(\overline{\mu}_p)=\sigma$, where $\sigma$ is the $q$-bar-core of both $\overline{\lambda}_p$ and $\overline{\mu}_p$. Then by Lemma \ref{lem4}(4) we have $\overline{(a\lambda)}_p=\overline{(b\mu)}_p=\sigma$, so using Lemma \ref{lem4}(1) we see that $\sigma$ is the $p$-bar-core and the $q$-bar-core of both $a\lambda$ and $b\mu$; in particular, $a\lambda$ and $b\mu$ have the same $p$-set. Moreover, by our assumption and the proof of the only `only if' part of the proposition above both $a\lambda$ and $b\mu$ have $q$-weighted $p$-quotient $\mathcal{Q}^q_p(\lambda)$. 

From the proof of \cite[Proposition 4.1]{F3} and the fact that $\sigma\in\overline{C}_p\cap\overline{C}_q$ it follows that for each $k\in\{0,\dots, q-1\}$, the elements $\Delta_j\sigma$ in the $p$-set (of $\sigma$, $a\lambda$ and $b\mu$) that are congruent to $k$ modulo $q$ form an arithmetic progression with common difference $q$. By the first paragraph of this proof we can therefore apply the level $q$ action to $a\lambda$ and arbitrarily reorder the elements $(a\lambda)^{(j)}\in\mathcal{Q}_p(a\lambda)$ such that $j\equiv k$ (mod $q$), for each $k$, without affecting the $q$-weighted $p$-quotient. 

Thus we can apply elements of $\mathfrak{W}_p$ to transform $a\lambda$ to $b\mu$, as $\mathcal{Q}^q_p(a\lambda)=\mathcal{Q}^q_p(b\mu)$ if and only if the multisets $[(a\lambda)^{(j)}|\Delta_j\sigma\equiv k\text{ (mod }q)]\text{ and }[(b\mu)^{(j)}|\Delta_j\sigma\equiv k\text{ (mod }q)]$ are equal for each $k\in\{0,\dots,q-1\}$, and it follows that $\lambda$ and $\mu$ lie in the same level $q$ orbit. 
\end{proof}

\section{Generalised bar-cores}
 
Now that we have covered all of the necessary definitions and basic results relating to the action of $\mathfrak{W}_p$, we arrive at the first of our main results. The following proposition is a generalisation of a theorem by Olsson \cite[Theorem 4]{O2} which states that the $q$-bar-core of a $p$-bar-core is again a $p$-bar-core, or in the notation used above, 
$$\overline{\text{wt}}_p(\lambda)=0\Rightarrow\overline{\text{wt}}_p(\overline{\lambda}_q)=0.$$ 

\begin{prop}\label{prop1}
For all bar partitions $\lambda$, 
$$\overline{\text{wt}}_p(\overline{\lambda}_q)\leq\overline{\text{wt}}_p(\lambda).$$
\end{prop}
\begin{proof}
We use induction on $\overline{\text{wt}}_q(\lambda)$, with the trivial case being that $\lambda$ is a $q$-bar-core. Assuming that this is not the case, we may find a removable $q$-bar: $y\in\lambda$ such that $y-q\not\in\mathcal{A}(\lambda)$. We will describe how to remove $q$-bars from $\lambda$ to obtain a new partition with the same $q$-bar-core as $\lambda$, with $q$-bar-weight strictly less than $\overline{\text{wt}}_q(\lambda)$, and with $p$-bar-weight no more than $\overline{\text{wt}}_p(\lambda)$.

Let $y$ be any part of $\lambda$ such that $y-q\not\in\mathcal{A}(\lambda)$. For any $x\in\mathcal{A}(\lambda)$ congruent to $y$ modulo $p$ such that $x-q\not\in\mathcal{A}(\lambda)$, replace $x$ with $x-q$, then replace $q-x\in\mathcal{A}(\lambda)$ with $-x$. We keep repeating this process until there are no more such $x$ (the process will terminate because $\lambda$ has finitely many removable $q$-bars), then we name our new bar partition $\nu$. Since each action corresponds to removing a $q$-bar from $\lambda$, and since we have removed at least one $q$-bar (replacing $y$ with $y-q$, and $q-y$ with $-y$, in $\mathcal{A}(\lambda)$), we have 
$$\overline{\nu}_q=\overline{\lambda}_q\text{ and }\overline{\text{wt}}_q(\nu)<\overline{\text{wt}}_q(\lambda).$$ 

We remarked earlier that the $p$-bar-weight of a bar partition $\lambda$ is equal to half the number of pairs $(x,a)\in\mathcal{A}(\lambda)\times\mathbb{N}$ such that $x-ap\not\in\mathcal{A}(\lambda)$. We will call such a pair $(x,a)$ a $p$\textbf{-bar-weight pair for} $\lambda$. It follows from our construction of $\nu$ that for any $x\not\equiv y,y-q,q-y,-y$ (mod $p$) and $a\in\mathbb{N}$, $(x,a)$ is a $p$-bar-weight pair for $\nu$ if and only if it is a $p$-bar-weight pair for $\lambda$. We will consider the remaining possibilities for the residue of $y$ modulo $p$ and show that in each case $\nu$ has no more $p$-bar-weight pairs than $\lambda$, hence $\overline{\text{wt}}_p(\nu)\leq\overline{\text{wt}}_p(\lambda)$. 

First suppose that $y\equiv q$ (mod $p$), so that we obtain $\mathcal{A}(\nu)$ by repeatedly replacing each $x\in\mathcal{A}(\lambda)$ such that $x\equiv q$ (mod $p$) and $q-x\in\mathcal{A}(\lambda)$ with $x-q$, then replacing $q-x$ with $-x$, until there are no more such $x$. Since in this case $y-q\equiv0\equiv q-y$ (mod $p$), we may compare the $p$-bar-weights of $\lambda$ and $\nu$ by counting how many of the three pairs $(x,a)$, $(x-q,a)$, $(x-2q,a)$ are $p$-bar-weight pairs for each of the two bar partitions when $x\equiv q$ (mod $p$) and $a\in\mathbb{N}$. We will do this by considering each of the four possibilities for the size of $X\coloneqq\mathcal{A}(\lambda)\cap\{x,x-q,x-2q\}$. 
\begin{flushleft}
$|X|=3$: If $x,x-q,x-2q\in\mathcal{A}(\lambda)$, then $x,x-q,x-2q\in\mathcal{A}(\nu)$, so the number of $p$-bar-weight pairs for $\lambda$, and for $\nu$, amongst the three pairs $(x,a)$, $(x-q,a)$, and $(x-2q,a)$ is equal to $3-|\mathcal{A}(\lambda)\cap\{x-ap,x-ap-q,x-ap-2q\}|$. 
\end{flushleft}
\begin{flushleft}
$|X|=2$: We have $\mathcal{A}(\nu)\cap\{x,x-q,x-2q\}=\{x-q,x-2q\}$, so clearly $(x,a)$ is not a $p$-bar-weight pair for $\nu$. If only one of the three pairs is a $p$-bar-weight pair for $\nu$, it must be $(x-q,a)$ as necessarily 
\begin{align*}
\mathcal{A}(\nu)\cap\{x-ap,x-ap-q,x-ap-2q\}=\{x-ap-2q\}.
\end{align*}
If $(x-q,a)$ and $(x-2q,a)$ are both $p$-bar-weight pairs for $\nu$, then we must have 
\begin{align*}
\mathcal{A}(\lambda)\cap\{x-ap,x-ap-q,x-ap-2q\}=\varnothing,\end{align*} 
so $\lambda$ also has two $p$-bar-weight pairs out of the three.
\end{flushleft}
\begin{flushleft}
$|X|=1$: We have $\mathcal{A}(\nu)\cap\{x,x-q,x-2q\}=\{x-2q\}$, so neither of $(x,a)$, $(x-q,a)$ can be $p$-bar-weight pairs for $\nu$. If $\mathcal{A}(\lambda)\cap\{x-ap,x-ap-q,x-ap-2q\}=\varnothing$, then exactly one of $(x,a)$, $(x-q,a)$, $(x-2q,a)$ is a $p$-bar-weight pair for $\lambda$. If 
\begin{align*}
|\mathcal{A}(\lambda)\cap\{x-ap,x-ap-q,x-ap-2q\}|\geq1,
\end{align*} 
then none of $(x,a)$, $(x-q,a)$, $(x-2q,a)$ can be $p$-bar-weight pairs for $\nu$ as 
\begin{align*}x-ap-2q\in\mathcal{A}(\nu). 
\end{align*}
\end{flushleft}
\begin{flushleft}
$|X|=0$: Since $\mathcal{A}(\nu)\cap\{x,x-q,x-2q\}=\varnothing$, none of $(x,a)$, $(x-q,a)$, $(x-2q,a)$ are $p$-bar-weight pairs for $\nu$. 
\end{flushleft}

Next suppose that $y\equiv0$ (mod $p$), so that we obtain $\mathcal{A}(\nu)$ by repeatedly replacing each $x\in\mathcal{A}(\lambda)$ such that $p|x$ and $q-x\in\mathcal{A}(\lambda)$ with $x-q$, then replacing $q-x$ with $-x$, until there are no much such $x$. Since $y\equiv-y$ (mod $p$), we can apply the same argument as above, when $y\equiv q$ (mod $p$), and conclude that $\nu$ has no more $p$-bar-weight pairs than $\lambda$ amongst $(x+q,a)$, $(x,a)$, $(x-q,a)$, and thus $\overline{\text{wt}}_p(\nu)\leq\overline{\text{wt}}_p(\lambda)$. 

Finally, suppose that $y\not\equiv q,0$ (mod $p$), so that $y\not\equiv-y$ and $y-q\not\equiv q-y$. In this case, we need only consider replacing all pairs $x,q-x\in\mathcal{A}(\lambda)$ such that $x\equiv y$ (mod $p$) with $x-q,-x$, so we are in a simpler situation; $\overline{\text{wt}}_p(\nu)\leq\overline{\text{wt}}_p(\lambda)$ since for any $x\equiv y$ (mod $p$) and $a\in\mathbb{N}$, $\nu$ has no more $p$-bar-weight pairs than $\lambda$ amongst $(x,a)$ and $(x-q,a)$.

Hence $\nu$ has no more $p$-bar-weight pairs than $\lambda$, and therefore has $p$-bar-weight no more than the $p$-bar-weight of $\lambda$. The result follows by induction.
\end{proof}

From this purely combinatorial result we obtain an interesting algebraic corollary. 

\begin{coro}For any $\mu\in\overline{\mathcal{P}}$, if $w$ is the weight of the $p$-block containing $[\mu]$, a spin representation of the symmetric group $\mathfrak{S}_r$ ($r\in\mathbb{N}$), and $[\lambda]$ is a spin representation of $\mathfrak{S}_{r+iq}$ corresponding to $\lambda\in\overline{\mathcal{P}}$ obtained by adding $q$-bars to $\mu$, for any $i\in\mathbb{N}$, then $[\lambda]$ belongs to a $p$-block of weight $\geq w$. 

In particular, if $[\mu]$ belongs to a $p$-block of weight $w>0$, then $[\lambda]$ belongs to a block of positive weight. 
\end{coro}

Next we will consider the set $\overline{C}_{p,q}$ containing all bar partitions $\lambda$ which satisfy 
$$\overline{\text{wt}}_p(\lambda)=\overline{\text{wt}}_p(\overline{\lambda}_q).$$

\begin{prop}\label{prop2}
For all $\lambda\in\overline{\mathcal{P}}$, the equality $\overline{\text{wt}}_p(\overline{\lambda}_q)=\overline{\text{wt}}_p(\lambda)$ holds if and only if there do not exist integers $a,b,c$ such that: 
\begin{align*}a\equiv b\:(\text{mod }p);\\
a\equiv c\:(\text{mod }q);\\
a,b+c-a\in\mathcal{A}(\lambda);\\
b,c\not\in\mathcal{A}(\lambda).\end{align*}
\end{prop}
\begin{proof}
Say that $(a,b,c)$ is a bad triple for $\lambda$ if $a,b,c$ satisfy the conditions above. When $(a,b,c)$ is bad, either $a>c$ or $b+c-a>b$; either way, since $a\equiv c$ (mod $q$) we find that $\lambda$ has a removable $q$-bar and is thus not a $q$-bar-core. Hence the proposition is true when $\lambda\in\overline{C}_q$. 

Now we assume $\lambda$ is not a $q$-bar-core, and choose $y\in\lambda$ such that $y-q\not\in\mathcal{A}(\lambda)$. We define a new bar partition $\nu$ as in the proof of Proposition \ref{prop1}: by repeatedly replacing pairs $x,q-x\in\mathcal{A}(\lambda)$ with $x-q$ and $-x$, respectively, when $x\equiv y$ (mod $p$). 

By induction it suffices to show that either:
\begin{enumerate}
\item  $\overline{\text{wt}}_p(\nu)=\overline{\text{wt}}_p(\lambda)$, and there is a bad triple for $\nu$ \textit{iff}. there is a bad triple for $\lambda$; or 
\item 
$\overline{\text{wt}}_p(\nu)<\overline{\text{wt}}_p(\lambda)$, and there is a bad triple for $\lambda$.
\end{enumerate}
Suppose first that there are no pairs $x,q-x$ satisfying 
\begin{align}\label{ass}x,q-x\not\in\mathcal{A}(\lambda)\text{ and }x\equiv y\text{ (mod }p\text{)}.\end{align} 
We first assume that $y\equiv q$ (mod $p$), and let $x\equiv q$ (mod $p$). Then there are eight different possibilities for the intersection of $\mathcal{A}(\lambda)$ and $\{x,x-q,x-2q\}$: 
\begin{align*}x,x-q,x-2q\in&\mathcal{A}(\lambda);\\
x,x-q\in&\mathcal{A}(\lambda)\not\ni x-2q;\\
x\in&\mathcal{A}(\lambda)\not\ni x-q,x-2q;\\
&\mathcal{A}(\lambda)\not\ni x,x-q,x-2q;\\
x,x-2q\in&\mathcal{A}(\lambda)\not\ni x-q;\\
x-q,x-2q\in&\mathcal{A}(\lambda)\not\ni x;\\
x-q\in&\mathcal{A}(\lambda)\not\ni x,x-2q;\\
x-2q\in&\mathcal{A}(\lambda)\not\ni x,x-q.\end{align*} 
However, the last four possibilities are all excluded by our assumption that there are no pairs $x,q-x$ satisfying (\ref{ass}), so we find that $\nu=\delta_0\lambda$. Therefore $\overline{\text{wt}}_p(\nu)=\overline{\text{wt}}_p(\lambda)$ by Lemma \ref{lem4}(3), and $(a,b,c)$ is bad for $\lambda$ exactly when $(\delta_0a,\delta_0b,\delta_0c)$ is bad for $\nu$, since $a\equiv b$ (mod $p$) $\Rightarrow \delta_0(b+c-a)=\delta_0b+\delta_0c-\delta_0a$.

If $y\equiv 0$ (mod $p$) we are in an identical situation to the above: $\nu=\delta_0\lambda$. 

When $y\not\equiv q,0$ (mod $p$), we have a similar situation: since there are no pairs $x,q-x$ satisfying (\ref{ass}), we have $\nu=\delta_i\lambda$, where $(i+1)q\equiv y$ (mod $p$). Hence $\overline{\text{wt}}_p(\nu)=\overline{\text{wt}}_p(\lambda)$, and $(a,b,c)$ is bad for $\lambda\Leftrightarrow (\delta_ia,\delta_ib,\delta_ic)$ is bad for $\nu$.

Finally, we assume that there is a pair $x,q-x$ satisfying (\ref{ass}), so that $(y,x,y-q)$ is a bad triple for $\lambda$. We argue that $\overline{\text{wt}}_p(\nu)<\overline{\text{wt}}_p(\lambda)$, as in the proof of Proposition \ref{prop1}: \newline
If $y\equiv q$ (mod $p$) and we let $z\coloneqq\max\{x,y\}$, $l\coloneqq\nicefrac{|x-y|}{p}$, then exactly one of $(z,l)$, $(z-q,l)$ is a $p$-bar-weight pair for $\lambda$ ($(z,l)$ if $x<y$, or $(z-q,l)$ if $x>y$). If $(z-2q,l)$ is a $p$-bar-weight pair for $\lambda$, then $(z-q,l)$ is a $p$-bar-weight pair for $\nu$ but neither of $(z,l)$, $(z-2q,l)$ are; and if $(z-2q,l)$ is not a $p$-bar-weight pair for $\lambda$, then none of $(z,l)$, $(z-q,l)$, $(z-2q,l)$ are $p$-bar-weight pairs for $\nu$. 

If instead we have $y\equiv 0$ (mod $p$), and again let $z\coloneqq$ max$\{x,y\}$ and $l\coloneqq\nicefrac{|x-y|}{p}$, then exactly one of $(z,l)$, $(z-q,l)$ is a $p$-bar-weight pair for $\lambda$. Now if $(z+q,l)$ is a $p$-bar-weight pair for $\lambda$, then $(z,l)$ is a $p$-bar-weight pair for $\nu$ but neither of $(z+q,l)$, $(z-q,l)$ are; and if $(z+q,l)$ is not a $p$-bar-weight pair for $\lambda$, then none of $(z+q,l)$, $(z,l)$, $(z-q,l)$ are $p$-bar-weight pairs for $\nu$. 

If $y\not\equiv q,0$ (mod $p$) then we define $z$ and $l$ as above so that exactly one of $(z,l),(z-q,l)$ is a $p$-bar-weight pair for $\lambda$ and neither is a $p$-bar-weight pair for $\nu$. 

Thus it follows from the proof of Proposition \ref{prop1} that there are less $p$-bar-weight pairs for $\nu$ then there are $p$-bar-weight pairs for $\lambda$. 
\end{proof}
\begin{thebibliography}{99}
\bibitem{BO} C. Bessenrodt and J. Olsson, Spin block inclusions, J. Algebra 306 (2006), 3--16.
\bibitem{F2} M. Fayers, Defect 2 spin blocks of symmetric groups and canonical basis coefficients, 2010, \url{http://www.maths.qmul.ac.uk/~mf/papers/spin_wt2.pdf}.
\bibitem{F3} M. Fayers, The n-bar-core of an m-bar-core, 2010, \url{http://www.maths.qmul.ac.uk/~mf/papers/barcores.pdf}.
\bibitem{F1} M. Fayers, A generalisation of core partitions, J. Comb. Theory, Ser. A 127 (2014), 58--84.
\bibitem{HH} P. Hoffman and J. Humphreys, Projective representations of the symmetric groups, Oxford Mathematical Monographs, Clarendon Press, 1992.
\bibitem{JK} G. James and A. Kerber, The representation theory of the symmetric group, Encyclop\ae dia of Mathematics and its Applications, vol. 16, Cambridge University Press, 1981.
\bibitem{MD} I. Macdonald, Symmetric functions and Hall polynomials 2nd edition, Oxford Mathematical Monographs, Oxford University Press, 1998.
\bibitem{O1} J. Olsson, Frobenius symbols for partitions and degrees of spin characters, Math. Scand. 61 (1987), 223--247.
\bibitem{O2} J. Olsson, A theorem on the cores of partitions, J. Comb. Theory, Ser. A 116 (2009), 733--740.
\end{thebibliography}
\end{document}
